% MT025 / explanation — canonical editable content.
题目只改了\textbf{数据损失}，没有加入权重正则。残差\(r=f(x)-y\)表示预测误差；正负表示高估或低估，绝对值表示错的大小。


用平方损失时，\(r=1,10,100\)分别贡献1、100、10,000。一个特别离谱的观测可能抵过很多普通样本，逼着回归线向它偏移。改用绝对损失后，三项是1、10、100，极端点仍受罚，但不会以平方速度扩大作用。


从更新方向看更直接。残差不为0时，
\[
\frac{d|r|}{dr}=\begin{cases}-1,&r<0,\\1,&r>0.\end{cases}
\]


而\(r^2\)的导数是\(2r\)。导数是损失随残差变化的局部斜率，斜率越大，一点误差变化带来的损失变化越强。绝对损失不会因为\(r\)从10变100就把斜率放大十倍，因此对大的输出异常值更稳健。\(r=0\)处有尖点，可以用次梯度等方法优化。


注意两条边界：损失本身仍随\(|r|\)增长，不能写成“大误差都罚一样”；这也不保证特征权重为0，因为绝对值没有加在\(w\)上。若\(x\)本身很大，参数梯度还含输入因子，对这类高杠杆点的稳健性要另看。


\textbf{出处：} 2021A期中 Q12，原卷第4页（10分）


\textbf{答案依据：} 2021A期中 Q12 随卷／配套参考答案；条件与解释经复核。
